\subsection{THETA3-A9: the exact factor-depth invariant and isotropic theta-frame cancellation}\label{A9}

There is a positive comparison behind the archive's divisibility
statements.  It uses more than complex harmonicity.  A specified
matrix of columns is Hermitian-isotropic in the completed CM
algebra, and has a unit minor.  The canonical theta-frame
correction restricts to an integral unit along those columns.
Its tensor power consequently does not change the derivative
image ideal.  This section proves that mechanism, identifies
the exact recovered statements, and distinguishes the
determinant exponent from the degree of the theta line bundle.

\paragraph{The three different integers.}
Let $i$ be the number of tangent rows in one determinant block,
let $a\geq0$ be its exponent, and let $k\geq0$ be the degree of
the theta line bundle.  On $i$ copies of a formal group $G$,
with an integral tangent frame $u_1,\ldots,u_i$, put
\[
 P_{i,a}(X)=\det(X_{rb})_{1\leq r,b\leq i}^{\,a}.
 \tag{A9a}
\]
The differential order is $ia$.  The number $k$ does not occur
in this polynomial.  For the normalized formal exponential
$\operatorname{Exp}_u$, define
\[
 \nu_{\det}(G,u,a)
   =\inf_{f\in R[[G^i]]}
       v_p\!\left(
         P_{i,a}(\partial)
            (f\circ\operatorname{Exp}_u)(0)\right),
 \qquad v_p(p)=1 .
 \tag{A9b}
\]
Here $R[[G^i]]$ means integral formal functions in an integral
coordinate system, and the exponential is evaluated separately
in each copy.  The notation is invariant under an integral
change of coordinates.  The operator is an actual derivative:
for $P=\sum_\alpha P_\alpha X^\alpha$ its value is
$\sum_\alpha P_\alpha\alpha![z^\alpha]f$.
It is not the divided Taylor coefficient.  These are exactly
\texttt{FGDepth.derivAt0}, \texttt{depth}, \texttt{detPow}, and
\texttt{detDepth} in
\texttt{Development/AbelianVarietyTheory/Definitions/Foundations/IsSatFrame.lean}.

In the archive, $\nu_{\mathrm{Fac}}$ is (A9b), expressed in an
integral chart and then made independent of the chart and of a
closed valued extension.  The canonical local depth
$\nu_{\mathrm{Loc}}(P,k)$ is instead the infimum of valuations
of
\[
 P(D)\left(g\exp\!\left(k\sum_{b=1}^n h_L(t_b)\right)\right)(0)
 \quad\text{over }g\in R[[G^n]],
 \tag{A9c}
\]
where $D$ denotes invariant differentiation and $h_L$ is the
canonical potential from A8.  Thus the exponential correction
is present in this definition.  Removing it without a further
argument would change the invariant.

\paragraph{A formal integral-unit lemma.}
Let $R$ be the valuation ring of a complete characteristic-zero
valued field $K$, with $v_p(p)=1$, and let $G$ be a commutative
formal group over $R$.  Let $\mathcal O$ be a commutative order, finite free
over $\Z$, acting on $G$, with involution $c\mapsto\bar c$.
Write $\ell$ for the normalized formal logarithm and
$\operatorname{Exp}$ for its inverse.  Suppose a formal series
$h(x,y)\in K[[G\times G]]$ vanishes on both axes and has the
following integral biextension cocycles:
\[
\begin{split}
 &\exp\bigl(h(x,z)+h(y,z)-h(x+_G y,z)\bigr),\\
 &\exp\bigl(h(z,x)+h(z,y)-h(z,x+_G y)\bigr),\\
 &\exp\bigl(h([c]x,y)-h(x,[\bar c]y)\bigr)
       \qquad(c\in\mathcal O)
\end{split}
 \tag{A9d}
\]
all belong to $1+(x,y,z)R[[x,y,z]]$, in the applicable variables.
No integrality of $h$ or $\exp h$ is assumed.
Fix $w_0\in\mathcal O$ and put
\[
 q(x)=\exp h(x,[w_0]x),\qquad
 \mathcal Q_n(x_1,\ldots,x_n)=\prod_{j=1}^n q(x_j).
 \tag{A9e}
\]
For the actual canonical potential, $w_0+\bar w_0=1$ and
$q=\exp h_L$.

Suppose matrices $a_N=(a_{N,bj})\in M_{r,n}(\mathcal O)$
satisfy
\[
 \sum_j\overline{a_{N,bj}}a_{N,b'j}\in p^N\mathcal O
       \quad\text{for every }b,b',
 \tag{A9f}
\]
and their formal group maps
\[
 \Phi_N:G^r\longrightarrow G^n,\qquad
 (\Phi_N)_j(x)=\sum_b^G[a_{N,bj}]x_b
\]
converge coefficientwise to a formal map $\Phi$.
Then $\Phi$ is integral and
\[
                 U:=\Phi^*\mathcal Q_n
                     \in 1+(x)R[[G^r]] .
 \tag{A9g}
\]
It suffices for convergence of $\Phi_N$ that the matrices
$d[a_{N,bj}]$ converge: the identity
$\ell([c]x)=d[c]\ell(x)$ expresses every finite jet of
$\Phi_N$ as a polynomial over $K$ in these matrix entries.

\textbf{Proof.}
Every $\Phi_N$ is integral.  Coefficientwise limits remain
integral because the valuation ring is closed.
For a fixed set of formal variables, define the additive
subgroup
\[
 \mathcal L
  =\{s\in(x)K[[x]]:\exp s\in1+(x)R[[x]]\}.
\]
It is a subgroup: multiplication of integral units corresponds
to addition, and the inverse of a series with integral
coefficients and constant term one is integral, by the
coefficient recursion for its inverse.
For integral formal points $x,y$, let $B(x,y)$ be the class of
$h(x,y)$ modulo $\mathcal L$.  Substitution of integral
zero-constant-term points into (A9d) shows that $B$ is
biadditive for the formal group law and satisfies
\[
             B([c]x,y)=B(x,[\bar c]y).
 \tag{A9h}
\]
In particular it is additive in the endomorphism applied to
either argument, including negative integer multiples.

Put $S_N=\sum_j h((\Phi_N)_j,[w_0](\Phi_N)_j)$.
Repeated use of biadditivity and (A9h) gives, modulo
$\mathcal L$,
\[
 S_N\equiv
  \sum_{b,b'}h\left(x_b,
       \left[w_0\sum_j\overline{a_{N,bj}}a_{N,b'j}\right]x_{b'}
            \right).
 \tag{A9i}
\]
Choose a fixed $\Z$-basis $e_1,\ldots,e_s$ of $\mathcal O$.
By (A9f), the endomorphism in brackets has the expression
\[
 \sum_\lambda p^N m_{N,bb'\lambda}e_\lambda,
 \qquad m_{N,bb'\lambda}\in\Z.
\]
Hence
\[
 S_N-R_N\in\mathcal L,\qquad
 R_N=\sum_{b,b',\lambda}
       p^N m_{N,bb'\lambda}h(x_b,[e_\lambda]x_{b'}).
 \tag{A9j}
\]
Every coefficient of $R_N$ tends to zero.  Indeed there are
only finitely many fixed series in (A9j), and integers have
nonnegative $p$-adic valuation.  All these series have zero
constant term.  Each coefficient of $\exp(-R_N)$ is therefore
a finite polynomial in coefficients of $R_N$, and
$\exp(-R_N)\to1$ coefficientwise.

By construction
$\exp(S_N-R_N)=\Phi_N^*\mathcal Q_n\exp(-R_N)$ is an integral
unit with constant term one.  Its coefficientwise limit is
$\Phi^*\mathcal Q_n$, because composition of zero-constant-term
formal maps uses only finitely many coefficients in each
degree.  Closedness of $R$ proves (A9g), and its constant term
is one.  This proves the lemma. \hfill$\square$

The proof uses the entire canonical frame, including its
higher terms.  It is distinct from the quadratic
mixed-Laplacian cancellation in A8.  Its additional input is
the Hermitian congruence (A9f).

\paragraph{Equality of derivative image ideals.}
Here is the precise consequence for arbitrary differential
polynomials.  Let $P$ be a polynomial in variables describing
chosen tangent directions on the $r$ source copies, and write
$E_u=\operatorname{Exp}_u$ for the corresponding formal
exponential map.  Suppose there is an integral zero-constant-term
map $\rho:G^n\to G^r$ with
\[
                   \rho\circ\Phi\circ E_u=E_u.
 \tag{A9k}
\]
Only this identity along the chosen exponential directions is
needed; a retraction on the entire formal group is sufficient
but is not required.  For any integer $k$, define
\[
 J_k(g)
   =P(\partial)\bigl(((\Phi^*g)\,U^k)
                         \circ E_u\bigr)(0),
          \qquad g\in R[[G^n]].
 \tag{A9l}
\]
Then
\[
 \{J_k(g):g\in R[[G^n]]\}
   =\{P(\partial)(f\circ E_u)(0):f\in R[[G^r]]\}.
 \tag{A9m}
\]
Both are $R$-submodules of $K$.  In particular their
valuation infima agree and are independent of $k$.

\textbf{Proof.}
For any integral $g$, the series
$(\Phi^*g)U^k$ is integral, because $U$ and its inverse are
integral.  This proves one inclusion.
Conversely take an integral $f$ and set
$g=\rho^*(fU^{-k})$.  This is integral.  Equation (A9k) implies
\[
 ((\Phi^*g)U^k)\circ E_u
   =((fU^{-k})\circ E_u)(U^k\circ E_u)
   =f\circ E_u.
 \]
Applying $P(\partial)$ proves the other inclusion.
The proof also establishes the assertion for negative $k$
as a local statement.  The global \texttt{IsNab} axioms in A8
still constrain actual sections only for nonnegative degree.
\hfill$\square$

For a product of determinant polynomials in disjoint blocks,
the ordinary image ideal is the product of the block image
ideals.  To see this without a tensor-completion assumption,
the operator has finite order and hence depends on a finite
jet of its argument.  Every monomial in that finite jet is a
product of block monomials, on which the functional factors.
This gives containment in the product ideal.  Conversely,
products of functions realizing block values realize the
products of those values.  The image ideals are finitely
generated by the finitely many jet monomials.  Their valuation
infima are minima, and valuation of a product adds.  Therefore
the total depth is the sum of the factor depths.

\paragraph{The exact recovered column condition.}
The preceding hypotheses are concrete in the recovered source.
In
\texttt{Development/AbelianVarietyTheory/Definitions/Divisors/NonemptyValExtOfOverLa.lean},
lines 327--344, \texttt{polyP} and \texttt{polyT} are products
of determinant powers
\[
 P_B(X)=\prod_f
   \det\left(\sum_{j=1}^n B_{f,bj}X_{\tau_f,rj}
                    \right)_{1\leq b,r\leq i_f}^{a_f}.
 \tag{A9n}
\]
For each $f$ with $a_f\ne0$, \texttt{IsCols} requires a
disjoint block of $2i_f$ section columns and
\texttt{IsIsoCols} requires the following three properties:
the matrix is supported in that block; an $i_f$ by $i_f$
minor is a unit; and for every $N$ it is congruent modulo
$p^N$, at the specified CM embedding, to a matrix
$x_N$ with entries in $\OO_E$ satisfying
$x_N\overline{x_N}^{\,t}\equiv0\pmod{p^N\OO_E}$.
The condition is stronger than belonging to \texttt{HHa}.
In particular a single nonzero column on an inert CM
elliptic factor does not satisfy it.

The source's prime-idempotent approximation constructs
limits of the $\OO_E$-actions whose tangent actions are the
prescribed entries $B_{f,bj}$ on the selected embedding
directions.  Applied blockwise, (A9f) persists because the
blocks are disjoint.  Its integral formal map is the
\texttt{InclSt} map $\psi$.
The unit minor constructs the \texttt{RetrSt} map $\rho$:
choose that minor, invert it in the local integral ring,
and use the limiting CM actions of its inverse entries.
The inverse entries remain in the same completed coefficient
ring by the adjugate formula and unit determinant.  On the
selected tangent directions the resulting matrices multiply
to the identity; uniqueness of the formal exponential then
gives (A9k).  Other blocks contribute zero in the selected
columns, by disjoint support.  This is precisely the
calculation in
\texttt{ColumnInclusionRetractionOfLimitSteps.lean}.
It does not require an integral projector onto one complex
embedding inside an inert local field.

Here are exact source locators for this mechanism.  In the
following list the common full prefix is
\texttt{Development/AbelianVarietyTheory/Proofs/DivisorProofs/}.
The definitions \texttt{IsIntFB}, \texttt{LogUnit},
\texttt{BiaddSt}, \texttt{LimFrameSt}, \texttt{InclSt}, and
\texttt{RetrSt} are in \texttt{MvIsIntMapMap.lean}.
The full proofs of the integral-unit mechanism are
\texttt{CanonGaugeBiadditiveOfLogSumActionLinear.lean},
\texttt{GaugeHermitianCombinationOfSteps.lean},
\texttt{LimitIsotropicFrameIntegralOfSteps.lean}, and
\texttt{ColumnInclusionIntegralOfLimitSteps.lean}.
The exact operator substitution is
\texttt{NablaOperatorFrameSeriesOfSteps.lean},
the inverse-unit construction is in
\texttt{ExistsIntegralGermOnCopiesSameValue.lean}, and
the depth equality is theorem \texttt{nuLocOf\_eq\_nuP}
in \texttt{ColumnPolynomialLocalDepthEqPlaceDepth.lean}.
It states, for every theta degree $k$ and every
\texttt{IsCols} matrix,
\[
 \nu_{\mathrm{Loc}}(P_B,k)
       =\sum_f\nu_{\mathrm{Fac}}(f,a_f),
 \tag{A9o}
\]
with the local valuation normalized by dividing the DVR
valuation by $v(p)$.  These proof bodies were read in this
run.  The source is the pinned secondary mirror commit
\texttt{7cbf7a0c32fca1badf05931879a37c3f8a2b03b4};
this is a source reconstruction, not a Lean build or a
transitive certification of all its dependencies.

\paragraph{The exact inert-prime bound and its scope.}
Write $j=\mathrm{Dw.q}$ and $P_0=\mathrm{Dw.p}$ to avoid
confusing the datum with the rational prime $p$.
The host contains the imaginary quadratic field generated
by its specified square root of $-d$, and
\texttt{InPsi} selects one of that field's embeddings.
For an odd prime unramified in $E$ and inert in this
quadratic field, the recovered lower bound is
\[
 \nu_{\mathrm{Fac}}(f,a)
 \geq
   \frac{(i-c_r)ap}{p^2-1}
       -i^2(\log_p a+1),\qquad
 0\leq c_r\leq i,\quad
 c_r\leq j\left(\frac32+\frac12\log_p a\right),
 \quad a\geq1,
 \tag{A9p}
\]
for a saturated rank-$i$ frame in an \texttt{InPsi} direction.
The precise formal-group theorem is
\texttt{detDepth\_lower} in
\texttt{DetDepthLowerBoundSaturatedEvenType.lean};
its application to the PEL datum is
\texttt{facBetween\_fine} in
\texttt{FactorDepthBoundedOddInertPrimeCmType.lean}.
Both full proof bodies were read, with their imported
formal-group radius and typed-basis inputs kept as source
inputs.  At an algebraic point the companion upper bound is
\[
 \nu_{\mathrm{Fac}}(f,a)
       \leq a\left(\frac{j}{p-1}
                    +\frac{(P_0-j)p}{p^2-1}\right);
 \tag{A9q}
\]
its exact statement is \texttt{nuFac\_le\_odd\_unram} in
\texttt{FactorDepthUpperBoundOddUnramifiedPrime.lean}.

For $j=0$ and $i=1$, (A9p) becomes
\[
 \nu_{\mathrm{Fac}}(f,a)
       \geq \frac{ap}{p^2-1}-\log_p a-1
       =\frac ap+\frac a{p^3}+\frac a{p^5}+\cdots
                                      -\log_p a-1.
 \tag{A9r}
\]
This is an exact recovered occurrence of the odd-power
divisibility slope in step 2 of the supplied timeline.
It is a statement about (A9b); applying it to a theta
section uses (A9o), its isotropic column data, and the
actual canonical normalization.
The timeline's period symbol must still be identified with
the chosen algebraic invariant differential before its
display can be substituted literally.

The source does transfer this depth to actual theta
generators.  Theorem
\texttt{harmonic\_value\_ge\_local\_depth} in
\texttt{HarmonicValueGeLocalDepthPrimeToP.lean}
uses a good integral chart, \texttt{IsNab}, an arithmetic
harmonic polynomial, a lower bound uniform in $k$, and
odd tame square-section data with level a unit at $p$.
It concludes the corresponding lower bound for the
canonical harmonic value.  Its proof reads reached
sections as integral local germs, uses the actual positive
exponential correction, and then applies the defining
infimum.  Thus the frame issue is resolved for these
specified operators and integral theta inputs.  It is
not resolved for every raw one-copy Taylor coefficient.

\paragraph{An actual elliptic instance and the comparison still required.}
C26 proves the integral-unit mechanism without relying on
the recovered archive.  For a good CM elliptic curve at an
inert prime, choose a unit $b$ in the completed CM order with
$b\bar b=-1$.  If $q(t)=\theta_{\rm red}(\lambda(t))/t$, then
\[
                 q(t)q([b]t)\in1+tR[[t]].
 \tag{A9s}
\]
The two columns $(1,b)$ have a unit minor and are
Hermitian-isotropic.  Actual normalized tame translated
theta sections on the two copies therefore restrict to an
integral function.  C19 gives the stronger explicit bound
\[
 v_p\!\left(D^a
    \bigl(s_{T_1}(t)s_{T_2}([b]t)\bigr)\big|_{t=0}\right)
   \geq
   \max\left\{0,\left\lfloor\frac{ap}{p^2-1}\right\rfloor-1\right\}.
 \tag{A9t}
\]
The differential is exactly $(D_1+bD_2)^a$ on the product.
This proves a positive-degree, actual-theta realization of
the strong slope on a specified formal isotropic direction.
The completed CM element need not define a map of complex
abelian varieties.

\paragraph{An explicit improvement in the product-formal case.}
There is a stronger quantitative statement under an additional
integral splitting hypothesis.  Suppose the selected $i$ tangent
directions on each of the $i$ source copies come from independent
one-dimensional strict height-two CM formal modules over a complete
DVR $R$, as in C19 and C21.  Normalize each invariant derivative
by its integral invariant differential.  If a CM isogeny is used
to obtain this product, require its degree to be prime to $p$;
it then gives an isomorphism of the formal completions.  Extra
formal coordinates may be retained as coefficient variables.
For $m\geq0$ put
\[
 \ell_p(m)=\max\left\{0,
          \left\lfloor\frac{mp}{p^2-1}\right\rfloor-1\right\},
 \qquad
 B_{i,p}(a)=
 \min_{\substack{\alpha\in\Z_{\geq0}^{i\times i}\!:\\
               \sum_b\alpha_{rb}=a\ (\forall r),\,
               \sum_r\alpha_{rb}=a\ (\forall b)}}
                   \sum_{r,b}\ell_p(\alpha_{rb}).
 \tag{A9u}
\]
The minimum is over a nonempty finite set, including when
$a=0$.  Then
\[
 \nu_{\det}(G,u,a)\geq B_{i,p}(a)
       \geq\max\left\{0,
                  \frac{iap}{p^2-1}-2i^2\right\}.
 \tag{A9v}
\]
This gives an explicit finite combinatorial certificate and a
constant-loss estimate in this product-formal regime.

\textbf{Proof.}
C19 proves operator divisibility, before evaluation at the
origin: the $m$th power of the invariant derivative on a strict
factor carries integral formal functions into
$p^{\ell_p(m)}R[[t]]$.  Apply this coefficientwise to all the
other variables.  The derivatives on the independent factors
commute, so the monomial $\prod_{r,b}D_{rb}^{\alpha_{rb}}$
carries integral functions into
$p^{\sum_{r,b}\ell_p(\alpha_{rb})}R[[G^i]]$.
Every monomial of $\det(D_{rb})^a$ has all row sums and all
column sums equal to $a$: each of its $a$ permutation factors
uses exactly one entry in every row and column.  Its coefficient
is an integer.  Consequently the determinant expansion introduces
no denominator, and every summand is divisible by
$p^{B_{i,p}(a)}$.  Evaluation and the defining infimum give
the first inequality in (A9v).  Finally
$\ell_p(m)\geq mp/(p^2-1)-2$ and $\ell_p(m)\geq0$.
Summing over the $i^2$ entries, whose total is $ia$, proves
the second inequality. \hfill$\square$

Under the isotropic-column and retraction hypotheses above,
(A9m) transfers the same bound to every theta degree $k$.
This improvement does not assert that a general typed formal
group in the recovered archive has such an integral product
decomposition.  For an actual 032 CM evaluation point, the
isogeny and differential normalizations must be checked at the
chosen prime before using it.

For the 032 top wedge, A8 identifies its actual complex
canonical derivative once the torus, section, frame, and
level are matched.  Its determinant in the original
section columns does not automatically satisfy
\texttt{IsCols}.  Replacing it by an isotropic doubled-column
operator changes the input unless a separate comparison
identifies the resulting automorphic coefficient.
This is now a concrete comparison of column matrices,
tensor degree, and derivative normalization.  C22 does
not contradict the recovered isotropic statement.
Finally, (A9s)--(A9t) are identities of formal germs and
their derivatives, not algebraic correspondences.
Within the requested Grothendieck framework, a passage
to homological motives must still specify the
correspondence and its polarization (cote44, scans 7--8)
and any Lefschetz inverse used for a degree change
(cote16, scans 43--46).
